\documentclass[12pt]{jlreq}
\usepackage{amsmath, amssymb}

\newcommand{\C}{\mathbb{C}}
\newcommand{\R}{\mathbb{R}}
\newcommand{\Z}{\mathbb{Z}}
\newcommand{\N}{\mathbb{N}}
\newcommand{\F}{\mathbb{F}}
\newcommand{\Prob}{\mathbb{P}}
\newcommand{\Exp}{\mathbb{E}}
\newcommand{\dd}{\,d}
\newcommand{\Ker}{\operatorname{Ker}}
\newcommand{\Impart}{\operatorname{Im}}
\newcommand{\Hom}{\operatorname{Hom}}

\begin{document}

\(0 < a < b\) に対し，
\[
  f_{a,b}(x) = \begin{cases} \frac{1}{x}, & a < |x| < b \text{ のとき}, \\ 0, & \text{その他のとき}, \end{cases}
\]
と定める．\(X = \{ g \in L^2(\R) \mid \|g\|_2 \le 1 \}\) とおくとき次の問いに答えよ．ただし，\(\|\cdot\|_2\) は \(L^2\) ノルムを表し，\(i = \sqrt{-1}\) とする．

(1) \(\sup_{g \in X} \|f_{a,b} * g\|_2 < \infty\) であることを示せ．ここで \(*\) は合成積 (convolution) を表す．

(2) \(\int_{-\infty}^\infty f_{a,b}(x) e^{-ix\xi} \dd x\) を，
  \[
    F(\xi) = \int_0^\xi \frac{\sin\eta}{\eta} \dd\eta
  \]
  を使って表せ．ここで \(\xi\) は実数である．

(3) \(\sup_{a,b} \sup_{g \in X} \|f_{a,b} * g\|_2 < \infty\) であることを示せ．ここで，\(\lim_{\xi \to \infty} F(\xi) = \frac{\pi}{2}\) は既知としてよい．

\end{document}